\subsection{Key-Pixel Feature Fusion}
    \label{subsec:kpff}
In echocardiography video segmentation, issues such as speckle noise, artifacts caused by signal attenuation, and low contrast between tissue and blood pose significant challenges to the feature extraction capabilities of models.
Linear models demonstrate rapid and stable memory capabilities in forming key-value associations. However, for inputs with rich spatial structures, such as images, relying solely on this type of associative memory may be insufficient for complex spatial reasoning. Sub-quadratic attention models like xLSTM~\cite{beck:24xlstm} widely utilize small-kernel 1D depth-wise separable convolutions to extract local semantic information. However, these spatial features often suffer from over-smoothing, leading to the loss of some global context.
\Cref{fig:kpff} shows that KPFF merges the convolution outputs learned from local key feature $\mathbf{F}_{\mathrm{K}}$, global key feature $\mathbf{F}_{\mathrm{Global}}$ and pixel-based feature $\mathbf{F}_{\mathrm{Pix}}$, each capturing information at different spatial scales, into a single network branch:
\begin{equation}
    \begin{aligned}
        \mathbf{F}_{\mathrm{Global}}        &= \mathrm{Expand} (\mathrm{GAP} (\mathbf{F}_{\mathrm{K}})), \\
        \mathbf{G}                          &= \sigma(\mathrm{Conv}_{\mathrm{gate}}(\mathbf{F}_{\mathrm{K}} + \mathbf{F}_{\mathrm{Global}})), \\
        \mathbf{F}_{\mathrm{fused}}         &= \mathbf{G} \cdot (\mathbf{F}_{\mathrm{K}} + \mathbf{F}_{\mathrm{Global}}) + (1 - \mathbf{G}) \cdot \mathbf{F}_{\mathrm{Pix}}, 
    \end{aligned}
    \label{eq:kpff}
\end{equation}
where $\sigma(\cdot)$ denotes the sigmoid function. 
When local features $\mathbf{F}_{\mathrm{K}}$ become unreliable due to speckle noise or signal dropout, $\mathbf{F}_{\mathrm{Global}}$ can provide a stable global context as a supplement. By fusing the smoothed global and local features with pixel-level features $\mathbf{F}_{\mathrm{Pix}}$, which retain high-frequency information, KPFF effectively preserves critical diagnostic details and prevents the loss of fine-grained information, such as valve opening and closing or regional wall motion abnormalities. 